\documentclass[12pt,a4paper]{article}
\usepackage[utf8]{inputenc}
\usepackage[portuguese]{babel}
\usepackage{amsmath}
\usepackage{amsfonts}
\usepackage{amssymb}
\usepackage{graphicx}
\usepackage{listings}
\usepackage{xcolor}
\usepackage{geometry}
\usepackage{hyperref}
\usepackage{titlesec}
\usepackage{booktabs}

\geometry{a4paper, left=30mm, right=20mm, top=30mm, bottom=20mm}

\definecolor{codegreen}{rgb}{0,0.6,0}
\definecolor{codegray}{rgb}{0.5,0.5,0.5}
\definecolor{codepurple}{rgb}{0.58,0,0.82}
\definecolor{backcolour}{rgb}{0.95,0.95,0.92}

\lstset{
    language=Python,
    backgroundcolor=\color{backcolour},   
    commentstyle=\color{codegreen},
    keywordstyle=\color{magenta},
    numberstyle=\tiny\color{codegray},
    stringstyle=\color{codepurple},
    basicstyle=\ttfamily\footnotesize,
    breakatwhitespace=false,         
    breaklines=true,                 
    captionpos=b,                    
    keepspaces=true,                 
    numbers=left,                    
    numbersep=5pt,                  
    showspaces=false,                
    showstringspaces=false,
    showtabs=false,                  
    tabsize=4
}

\titleformat{\section}{\normalfont\Large\bfseries\color{blue!70!black}}{\thesection}{1em}{}
\titleformat{\subsection}{\normalfont\large\bfseries\color{blue!50!black}}{\thesubsection}{1em}{}

\title{\textbf{Documentação Técnica e Especificação do Compilador de Expressões Matemáticas}}
\author{Desenvolvimento de Compiladores | Relatório Técnico}
\date{\today}

\begin{document}

\maketitle

\begin{abstract}
Este documento apresenta a análise técnica, modelagem matemática e arquitetura do compilador/interpretador de expressões matemáticas desenvolvido em Python. O sistema realiza as etapas fundamentais de análise léxica (através de expressões regulares) e análise sintática preditiva descendente recursiva (\textit{Recursive Descent Parsing}) com avaliação de expressões e persistência de variáveis em tempo de execução.
\end{abstract}

\tableofcontents
\newpage

\section{Introdução e Arquitetura Global}
A engenharia de compiladores estuda as técnicas necessárias para traduzir ou interpretar linguagens de programação de alto nível. O código analisado propõe um interpretador de expressões aritméticas baseado em uma Gramática Livre de Contexto (GLC) com prioridades de operadores e associatividade explícitas.

O fluxo de dados do compilador segue uma estrutura linear dividida em:
\begin{enumerate}
    \item \textbf{Cadeia de Entrada:} Código fonte bruto fornecido pelo usuário.
    \item \textbf{Analisador Léxico (Tokenizer):} Quebra a cadeia de texto em tokens tipados baseados em expressões regulares.
    \item \textbf{Analisador Sintático (Parser):} Consome os tokens baseando-se em regras gramaticais e calcula de forma imediata o resultado semântico.
\end{enumerate}

\section{Análise Léxica (Lexer)}
O processo de tokenização varre a string de entrada comparando-a com uma lista de padrões descritos via expressões regulares (Regex). 

\subsection{Tabela de Especificação de Tokens}
Abaixo são mapeados os tipos de tokens aceitos pelo motor do \textit{Lexer}:

\begin{table}[h!]
\centering
\caption{Especificação de Expressões Regulares dos Tokens}
\begin{tabular}{lll}
\toprule
\textbf{Token Class} & \textbf{Expressão Regular} & \textbf{Descrição} \\
\midrule
\texttt{NUMBER} & \verb|\d+(\.\d+)?| & Números inteiros ou de ponto flutuante \\
\texttt{IDENT} & \verb|[A-Za-z_][A-Za-z0-9_]*| & Variáveis e Identificadores \\
\texttt{POW} & \verb|\*\*|\verb|^\^| & Operadores de Exponenciação \\
\texttt{PLUS} / \texttt{MINUS} & \verb|\+| / \verb|-| & Adição e Subtração \\
\texttt{MUL} / \texttt{DIV} & \verb|\*| / \verb|/| & Multiplicação e Divisão \\
\texttt{ASSIGN} & \verb|=| & Atribuição de variáveis \\
\texttt{LPAREN} / \texttt{RPAREN} & \verb|\(| / \verb|\)| & Delimitadores de Parênteses \\
\texttt{LBRACK} / \texttt{RBRACK} & \verb|\[| / \verb|\]| & Delimitadores de Colchetes \\
\texttt{LBRACE} / \texttt{RBRACE} & \verb|\{| / \verb|\}| & Delimitadores de Chaves \\
\bottomrule
\end{tabular}
\end{table}

\section{Análise Sintática e Gramática Formal}
O \textit{Parser} utiliza uma abordagem Descendente Recursiva (\textit{Recursive Descent Parser}), onde cada nível de precedência matemática é isolado em um método específico (\texttt{expr}, \texttt{term}, \texttt{power}, \texttt{factor}).

\subsection{Gramática na Notação EBNF}
A gramática livre de contexto pode ser descrita formalmente como:

\begin{align*}
\mathit{Parse} &\rightarrow \mathit{IDENT} \ \mathbf{=} \ \mathit{Expr} \ \mid \ \mathit{Expr} \\
\mathit{Expr} &\rightarrow \mathit{Term} \ \{ (\mathbf{+} \mid \mathbf{-}) \ \mathit{Term} \} \\
\mathit{Term} &\rightarrow \mathit{Power} \ \{ (\mathbf{*} \mid \mathbf{/}) \ \mathit{Power} \} \\
\mathit{Power} &\rightarrow \mathit{Factor} \ [ \mathbf{**} \mid \mathbf{\wedge} \ \mathit{Power} ] \\
\mathit{Factor} &\rightarrow \mathit{NUMBER} \mid \mathit{IDENT} \mid \mathbf{(}\mathit{Expr}\mathbf{)} \mid \mathbf{[}\mathit{Expr}\mathbf{]} \mid \mathbf{\{}\mathit{Expr}\mathbf{\}}
\end{align*}

A regra $\mathit{Power}$ introduz a associatividade à direita, permitindo resolver corretamente cálculos encadeados como $2 \wedge 3 \wedge 2 = 2^{(3^2)} = 512$.

\section{Implementação em Python}
Abaixo encontra-se o código fonte unificado que serve como fundação estrutural deste sistema compilador:

\begin{lstlisting}
import re

ENV = {}

def compiler(code):
    try:
        tokens = tokenize(code)
        parser = Parser(tokens, ENV)
        return parser.parse()
    except Exception as e:
        return f"Erro: {e}"

TOKEN_SPECIFICATION = [
    ('NUMBER',   r'\d+(\.\d+)?'),
    ('IDENT',    r'[A-Za-z_][A-Za-z0-9_]*'),
    ('POW',      r'\*\*|\^'),
    ('PLUS',     r'\+'),
    ('MINUS',    r'-'),
    ('MUL',      r'\*'),
    ('DIV',      r'/'),
    ('ASSIGN',   r'='),
    ('LPAREN',   r'\('),
    ('RPAREN',   r'\)'),
    ('LBRACK',   r'\['),
    ('RBRACK',   r'\]'),
    ('LBRACE',   r'\{'),
    ('RBRACE',   r'\}'),
    ('SKIP',     r'[ \t\n]+'),
    ('MISMATCH', r'.'),
]

def tokenize(code):
    tok_regex = '|'.join(f'(?P<{name}>{regex})' for name, regex in TOKEN_SPECIFICATION)
    tokens = []
    for mo in re.finditer(tok_regex, code):
        kind = mo.lastgroup
        value = mo.group()
        if kind == 'NUMBER':
            tokens.append(('NUMBER', float(value)))
        elif kind == 'IDENT':
            tokens.append(('IDENT', value))
        elif kind == 'SKIP':
            continue
        elif kind == 'MISMATCH':
            raise SyntaxError(f'Caractere inesperado: {value}')
        else:
            tokens.append((kind, value))
    tokens.append(('EOF', None))
    return tokens

class Parser:
    def __init__(self, tokens, env=None):
        self.tokens = tokens
        self.pos = 0
        self.current_token = self.tokens[self.pos]
        self.env = env if env is not None else {}

    def error(self, message="Erro de sintaxe"):
        raise SyntaxError(message)

    def advance(self):
        self.pos += 1
        if self.pos < len(self.tokens):
            self.current_token = self.tokens[self.pos]

    def eat(self, token_type):
        if self.current_token[0] == token_type:
            self.advance()
        else:
            self.error(f"Esperado {token_type}, encontrado {self.current_token[0]}")

    def factor(self):
        token = self.current_token
        if token[0] == 'NUMBER':
            self.eat('NUMBER')
            return token[1]
        elif token[0] == 'IDENT':
            name = token[1]
            self.eat('IDENT')
            if name in self.env:
                return self.env[name]
            self.error(f"Variavel nao definida: {name}")
        elif token[0] == 'LPAREN':
            self.eat('LPAREN')
            result = self.expr()
            self.eat('RPAREN')
            return result
        elif token[0] == 'LBRACK':
            self.eat('LBRACK')
            result = self.expr()
            self.eat('RBRACK')
            return result
        elif token[0] == 'LBRACE':
            self.eat('LBRACE')
            result = self.expr()
            self.eat('RBRACE')
            return result
        self.error(f"Fator invalido: {token[1]}")

    def term(self):
        node = self.power()
        while self.current_token[0] in ('MUL', 'DIV'):
            op = self.current_token[0]
            self.advance()
            if op == 'MUL':
                node = node * self.power()
            elif op == 'DIV':
                denom = self.power()
                if denom == 0:
                    raise ZeroDivisionError("Divisao por zero.")
                node = node / denom
        return node

    def power(self):
        if self.current_token[0] in ('NUMBER', 'IDENT', 'LPAREN', 'LBRACK', 'LBRACE'):
            node = self.factor()
            if self.current_token[0] == 'POW':
                self.eat('POW')
                exponent = self.power()
                node = node ** exponent
            return node
        self.error(f"Fator invalido para potencia: {self.current_token}")

    def expr(self):
        node = self.term()
        while self.current_token[0] in ('PLUS', 'MINUS'):
            op = self.current_token[0]
            self.advance()
            if op == 'PLUS':
                node = node + self.term()
            elif op == 'MINUS':
                node = node - self.term()
        return node

    def parse(self):
        if self.current_token[0] == 'IDENT' and self.pos + 1 < len(self.tokens) and self.tokens[self.pos + 1][0] == 'ASSIGN':
            name = self.current_token[1]
            self.eat('IDENT')
            self.eat('ASSIGN')
            value = self.expr()
            self.env[name] = value
            if self.current_token[0] != 'EOF':
                self.error("Tokens extras detectados apos atribuicao.")
            return value

        result = self.expr()
        if self.current_token[0] != 'EOF':
            self.error("Tokens extras detectados no fim da expressao.")
        return result
\end{lstlisting}

\section{Conclusão}
A arquitetura proposta demonstra viabilidade para execução confiável de expressões matemáticas aninhadas. O uso do padrão de análise descendente recursiva aliada ao encapsulamento léxico por expressões regulares confere modularidade e extensibilidade para futuras implementações, como funções nativas (\texttt{sin}, \texttt{cos}, \texttt{log}) ou blocos de controle estruturados.

\end{document}
